\hnsection{李代数理论}

一旦建立了变换群与李代数之间的对应，理论便转向了更加代数化的方向，并以深入研究李代数为中心。†

1888 年至 1894 年是第一个短暂时期；这一时期以 Engel 及其学生 Umlauf，尤其是 Killing 和 E. Cartan 的工作为标志，在复李代数方面取得了一系列惊人的成果。我们此前已经看到，可解李代数的概念源于 Lie 本人；他曾证明（在复数情形下）可解线性李代数化为三角形形式的定理 ({[}4{]}, vol.~1, p.~270)。\(^{\ddagger}\) Killing 注意到 {[}11{]}，李代数中存在最大的可解理想（我们现在称之为根基），且该李代数对其根基的商的根基为零；他把根基为零的李代数称为半单李代数，并证明它们是单代数的乘积（后一个概念已经由 Lie 引入，他曾证明“经典”代数的单纯性 ({[}4{]}, vol.~3, p.~682)）。

另一方面，Killing 在李代数中引入了特征方程 \(\det(\mathrm{ad}(x)-\omega.1)=0\) ，Lie 在研究包含给定李代数元素的二维李子代数时已经遇到过这个方程。我们将在本书其他《历史注记》中回到对这些方法的分析：Killing 深入研究半单代数的“通用”特征方程的根的性质时，取得了他最卓越的成果，即完全确定（复）单李代数。§

Killing 证明，可解代数的导出代数是“秩为 0”的（这意味着对代数的每个元素 x，ad x 都是幂零的）。几乎与此同时，Engel 证明了“秩为 0”的代数是可解的（这个命题实质上就是我们在第一章 § 4，第 2 号中称为 Engel 定理的命题）。另一方面，E. Cartan 在其学位论文中引入了我们现在称为“Killing 型”的对象，并建立了两个基本判据，它们借助该型刻画了可解李代数和半单李代数。

Killing 曾断言 ({[}11{]}, IV)，李代数的导出代数是一个半单代数与其根基之和，而该根基是幂零的，但他的证明并不完整。不久以后，E. Cartan 未经证明便宣布 ({[}12{]}, vol.~I₁, p.~104)，更一般地说，每个李代数都是其根基与一个半单子代数之和；这一时期在这个方向上无可争议地确立的唯一结果，是 Engel 的一个定理，它断言每个非可解李代数中都存在一个 3 维单李子代数。Cartan 这一断言的第一个发表的证明（针对复李代数）归于 E. E. Levi {[}18{]}；J. H. C. Whitehead 于 1936 年给出了另一个证明（在实数情形下同样有效）{[}26 a{]}。1942 年，A. Malcev 通过关于“Levi 截面”在共轭意义下唯一性的定理完成了这一结果。

Lie 从其最初的著作起，就一直关注任意李代数与线性李代数之间的同构问题。他曾相信，通过考察伴随表示，他已经肯定地解决了这个问题（并由此推导出他的“第三定理”的一个证明）\([3\ d]\) ；后来他意识到自己的证明只对中心为零的李代数正确，又给出了另一个错误的证明（ \([3\ f]\) ，§ 3, no. 9），随后认识到（ \([3\ g]\) ，p.~231）这个问题仍未解决。事实上，这个问题长期悬而未决，直到 1935 年 Ado 才肯定地解决了它 {[}27{]}。另一方面，Lie 实质上提出了确定最小维数的单李代数的线性表示这一问题，并已对经典群解决了它；Cartan 在其学位论文中也解决了例外单代数的这一问题 \(^{\dagger}\) ；二十年后，他所采用的方法被他推广，用以得到实或复单李代数的所有不可约表示。

线性表示的完全可约性这一性质，似乎是 Study 首先（以几何形式）遇到的。在一篇未发表的手稿中（但已在 {[}4{]}，vol.~3, pp.~785--788 中引述），他证明了李代数 \(\mathbf{SL}(2,\mathbf{C})\) 的线性表示具有此性质，并对 \(\mathbf{SL}(3,\mathbf{C})\) 和 \(\mathbf{SL}(4,\mathbf{C})\) 取得了部分结果。Lie 和 Engel 借此猜想，\(\mathbf{SL}(n,\mathbf{C})\) 的完全可约性定理对所有 \(n\) 成立。半单李代数线性表示的完全可约性由 H. Weyl 于 1925 年† 通过一种全局型论证（见后文）建立。第一个代数证明由 Casimir 和 van der Waerden 于 1935 年取得 {[}32{]}；此后 R. Brauer {[}31{]}（这就是我们复现的证明）和 J. H. C. Whitehead {[}26 b{]} 又给出了其他代数证明。

最后，在研究指数映射的过程中（cf.~infra），H. Poincaré ({[}14{]}, vol.~3) 考察了由李代数各算子生成的全阶微分算子的结合代数；他实质上证明，若 \((\mathbf{X}_i)_{1\leqslant i\leqslant n}\) 是该李代数的一个基，则由 \(\mathbf{X}_i\) 生成的结合代数以 \(\mathbf{X}_i\) 的某些对称函数为基（这些函数是将给定单项式的因子作所有置换后得到的非交换“单项式”之和）。他的证明的本质是代数性的，这使他得以得到我们在第一章中抽象定义的包络代数结构。类似的证明由 G. Birkhoff {[}29 b{]} 和 E. Witt {[}30{]} 于 1937 年给出。

上述大多数结果限于实或复李代数，只有它们按通常意义对应于李群。Jacobson 开始研究定义在 \(\mathbf{R}\) 或 \(\mathbf{C}\) 以外的域上的李代数 {[}28 a{]}；他证明经典结果的大部分（即第一章的结果）对任意特征为零的域仍然成立。

